% Evans 8.4: printed pp.458--471 / English PDF pp.475--488.
% Chinese: printed pp.323--332 / scan pp.341--350.
\Needspace{55mm}
\section{Constraints}
\label{sec:8-4}
\Needspace{45mm}
\subsection{Integral constraints and a scalar multiplier}
On connected $U$, minimize $I[u]=\frac12\int_U|Du|^2$ over
\[
\mathcal A=\{u\in H^1_0(U):J[u]=0\},\qquad J[u]=\int_U G(u).
\]
Assume $g=G'$ satisfies $|g(z)|\leq C(1+|z|)$.
Then $|G(z)|\leq C(1+z^2)$, so the constraint is well defined.
\Needspace{55mm}
\begin{theorem}[Constrained existence; source \S8.4.1, Theorem 1]
If $\mathcal A\ne\varnothing$, it contains an energy minimizer.
\end{theorem}
\begin{expandedproof}[compactness preserves the nonlinear integral]
An energy-bounded minimizing sequence is bounded in $H^1_0$.
After extraction $u_k\rightharpoonup u$ in $H^1_0$ and $u_k\to u$ in $L^2$.
The mean value theorem and Cauchy--Schwarz give
\[
 |J[u_k]-J[u]|\leq C\|u_k-u\|_2
 (|U|^{1/2}+\|u_k\|_2+\|u\|_2)\longrightarrow0.
\]
Thus $u$ is feasible, and weak lower semicontinuity of the Dirichlet norm
makes it minimizing. Every member of this particular class has finite energy.
\end{expandedproof}
\Needspace{55mm}
\begin{theorem}[Lagrange multiplier; source Theorem 2]
For a minimizer $u$ there is $\lambda\in\R$ such that
\begin{equation}\label{eq:8-multiplier}
 \int_U Du\cdot D\phi=\lambda\int_Ug(u)\phi
 \quad(\phi\in H^1_0(U)).
\end{equation}
\end{theorem}
\begin{expandedproof}[an admissible curve and the degenerate case]
First suppose $g(u)\not\equiv0$. By density choose $\psi\in H^1_0$ with
$b=\int_Ug(u)\psi\ne0$. For any $\phi\in H^1_0$, set
$j(s,t)=J[u+s\phi+t\psi]$. Its derivatives are continuous by the
linear growth bound on $g$ and dominated convergence on this finite
parameter plane. Since $j_t(0,0)=b$, the ordinary implicit function
theorem supplies $t=\eta(s)$ with $j(s,\eta(s))=0$,
$\eta(0)=0$, and $\eta'(0)=-\int g(u)\phi/b$.
Minimality on this curve gives
$\int Du\cdot D\phi+\eta'(0)\int Du\cdot D\psi=0$.
Thus \eqref{eq:8-multiplier} holds with
$\lambda=(\int Du\cdot D\psi)/b$, independent of $\phi$.

If $g(u)=0$ a.e., the chain rule with truncation gives
$G(u)\in W^{1,1}$ and $D(G(u))=g(u)Du=0$.
The product is integrable by Cauchy--Schwarz. Connectedness makes $G(u)$
constant, and $J[u]=0$ makes it zero. To identify its boundary value,
approximate $u$ in $H^1_0$ by compactly supported smooth functions:
the linear growth of $g$ implies convergence of their compositions in
$W^{1,1}$, so $G(u)-G(0)$ has zero trace. Hence $G(0)=0$ and zero is
feasible. Minimality gives $I[u]=0$, whence $u=0$. Both sides of
\eqref{eq:8-multiplier} vanish for every $\lambda$.
This final conclusion is special to this energy and constraint, not a
universal multiplier theorem without a constraint qualification.
\end{expandedproof}
\Needspace{55mm}
\PDEAnnotation{退化约束不能直接除法}{只有 $J'[u]\ne0$ 时才能用隐函数定理并除以 $\int g(u)\psi$。本例退化时另证 $u=0$；一般约束问题不保证退化点存在普通 Lagrange 乘子。}

\Needspace{45mm}
\subsection{The obstacle problem and a variational inequality}
Let $f\in L^2(U)$, let $h$ be a smooth obstacle, and assume
$K=\{w\in H^1_0(U):w\geq h\text{ a.e.}\}\ne\varnothing$.
Minimize $I[w]=\frac12\|Dw\|_2^2-(f,w)$.
\Needspace{55mm}
\begin{theorem}[Obstacle minimizer; source \S8.4.2, Theorems 3--4]
There is a unique minimizer $u\in K$, characterized equivalently by
\begin{equation}\label{eq:8-VI}
 \int_U Du\cdot D(w-u)\geq\int_Uf(w-u)\qquad(w\in K).
\end{equation}
\end{theorem}
\begin{expandedproof}[one-sided variations and the converse]
Young and Poincar\'e imply
$I[w]\geq\frac14\|Dw\|_2^2-C\|f\|_2^2$.
Weak limits of minimizing sequences stay in $K$: strong $L^2$ convergence
has an a.e. convergent subsequence preserving $w\geq h$.
The direct method gives existence. The midpoint identity gives strict
convexity on $H^1_0$, hence uniqueness.
For $w\in K$, the segment $u+s(w-u)\in K$ for $0\leq s\leq1$.
Its right energy derivative at zero is nonnegative, giving \eqref{eq:8-VI}.
Conversely, expand the quadratic energy exactly:
\[
 I[w]-I[u]=\int_U Du\cdot D(w-u)-f(w-u)
                    +\tfrac12\|D(w-u)\|_2^2\geq0.
\]
This proves equivalence, including sufficiency of the inequality.
\end{expandedproof}
For every nonnegative $\phi\in C_c^\infty(U)$, $u+\phi\in K$ gives
$\mu=-\Delta u-f\geq0$ in distributions. If $u$ is continuous, then
$O=\{u>h\}$ is open; on a compact subset of $O$ both signs of a small
variation are allowed, so $-\Delta u=f$ there. If also $u\in H^2_{\rm loc}$,
then $\mu$ is a function and
\[
 u-h\geq0,\quad\mu\geq0,\quad (u-h)\mu=0\quad\text{a.e.}
\]
Interior continuity and $C^{1,1}_{\rm loc}$ regularity for smooth obstacle
and forcing require external obstacle regularity theory; they are not
proved by the direct method above. No global $W^{2,\infty}$ estimate
is claimed without specifying its boundary compatibility hypotheses.
The free boundary is $\partial\{u>h\}\cap U$.
\Needspace{55mm}
\PDEAnnotation{接触集与反力}{约束只允许向上变化，故反力 $\mu=-\Delta u-f$ 非负。在非接触区可双向变化，反力为零。可行集、反力和互补条件要使用同一符号约定。}

\PDEBodyDiagram{8.3}{An exact obstacle solution with two free-boundary points.}
{On $(-1,1)$ take $f=0$ and $h(x)=1-4x^2/3$. The minimizer is $u=h$ on $|x|\leq1/2$, and $u=4(1-|x|)/3$ outside. Values and first derivatives match at $\pm1/2$, so $-u''=8/3$ on the contact interval and zero elsewhere, with no point masses.}
{ch08-obstacle}

\Needspace{45mm}
\subsection{Pointwise constraints: harmonic maps}
Consider
$\mathcal A=\{u\in H^1(U;\R^m):u-g\in H^1_0,\ |u|=1\text{ a.e.}\}$,
and $I[u]=\frac12\int_U|Du|^2$.
A minimizing sequence is bounded in $H^1$ ($|u|=1$ fixes its $L^2$ norm).
A strong $L^2$ subsequence preserves the sphere constraint a.e.; the trace
condition is weakly closed, so a minimizer exists if the class is nonempty.
\Needspace{55mm}
\begin{theorem}[Harmonic map equation; source \S8.4.3, Theorem 5]
A minimizer satisfies
\begin{equation}\label{eq:8-harmonic-map}
 \int_U Du:D\phi=\int_U|Du|^2u\cdot\phi
 \quad(\phi\in H^1_0(U;\R^m)\cap L^\infty(U;\R^m)).
\end{equation}
\end{theorem}
\begin{expandedproof}[normalization generates tangent variations]
For bounded $\phi$, $u_s=(u+s\phi)/|u+s\phi|$ is well defined for
$|s|\|\phi\|_\infty<1/2$. Sobolev composition and traces show
$u_s\in\mathcal A$. Its derivative at zero in $H^1$ is
$\dot u_0=\phi-(u\cdot\phi)u$.
To justify this derivative, derivatives of the normalization map are
uniformly bounded on $|u+s\phi|\geq1/2$; its spatial derivative and
parameter derivative are dominated in $L^2$ by
$C(|D\phi|+|Du|\,\|\phi\|_\infty)$.
Thus differentiating the energy is legitimate.
Since $u\cdot\partial_i u=0$ a.e.,
\[
 Du:D((u\cdot\phi)u)=|Du|^2(u\cdot\phi).
\]
The vanishing derivative is precisely \eqref{eq:8-harmonic-map}.
The multiplier is the function $|Du|^2\in L^1$, explaining why bounded
variations appear in the weak formulation.
\end{expandedproof}
\Needspace{55mm}
\PDEAnnotation{约束曲率进入方程}{球面约束的切向投影是 $\phi-(u\cdot\phi)u$。对 $|u|^2=1$ 求导给 $u\cdot u_i=0$；再次求导产生 $|Du|^2u$，这里乘子是函数而非一个常数。}

\Needspace{45mm}
\subsection{Divergence constraints and pressure}
For $f\in L^2(U;\R^n)$ minimize
$I[v]=\frac12\int|Dv|^2-\int f\cdot v$ on the closed linear space
$V_\sigma=\{v\in H^1_0(U;\R^n):\operatorname{div}v=0\}$.
Coercivity and strict convexity give a unique $u$, with
$\ell(v)=\int Du:Dv-f\cdot v=0$ for $v\in V_\sigma$.
To recover pressure we use the following external divergence theorem:
on a bounded connected smooth domain there is a bounded right inverse
\[
 \mathcal B:L^2_0(U)\to H^1_0(U;\R^n),\quad
 \operatorname{div}\mathcal Bq=q,\quad\|\mathcal Bq\|_{H^1}\leq C\|q\|_2.
\]
Here $L^2_0$ denotes zero mean. This is the divergence-solving estimate
also invoked without proof in the source, often realized by a Bogovskii
operator. Its construction is not proved here.
\Needspace{55mm}
\begin{theorem}[Pressure multiplier; source \S8.4.4, Theorem 6]
There is a unique mean-zero $p\in L^2(U)$ satisfying
\[
 \int_U Du:Dv=\int_U f\cdot v+p\operatorname{div}v
 \quad(v\in H^1_0(U;\R^n)).
\]
Thus $-\Delta u+Dp=f$, $\operatorname{div}u=0$ in distributions.
The global conclusion uses the external right inverse just stated; the
source proves a local pressure statement.
\end{theorem}
\begin{expandedproof}[factoring a functional through divergence]
The functional $q\mapsto\ell(\mathcal Bq)$ is bounded on $L^2_0$.
Riesz representation gives mean-zero $p$ with
$\ell(\mathcal Bq)=\int pq$. For any $v\in H^1_0$,
$q=\operatorname{div}v$ has zero mean and
$v-\mathcal Bq\in V_\sigma$. Hence
$\ell(v)=\ell(\mathcal Bq)=\int p\operatorname{div}v$.
Surjectivity of divergence gives uniqueness in $L^2_0$ and
$\|p\|_2\leq C(\|Du\|_2+\|f\|_2)$.
Without mean normalization constants are undetermined, since
$\int\operatorname{div}v=0$. This argument needs connectedness and the
right inverse, not a claim that every curl-free field is globally a gradient.
\end{expandedproof}
\Needspace{55mm}
\PDEAnnotation{压力只确定到常数}{测试速度具有零边界，故散度积分为零；压力加常数不改变弱式。右逆估计将消失在无散度空间上的泛函表示为 $p\operatorname{div}v$，并给出压力范数控制。}

\Needspace{45mm}
\subsection{Nonlinear incompressibility}
In three-dimensional elasticity, $u$ is the deformation map (the new
position), not the displacement $u-x$. Local incompressibility is
$\det Du=1$. For an actual displacement $v=u-x$ it is
$\det(I+Dv)=1$, whose first-order linearization is $\operatorname{div}v=0$.
\Needspace{55mm}
\begin{theorem}[Determinant constraint; source Theorem 7]
For $q>3$, let $L(P,x)$ be smooth, convex in $P$, and satisfy
\eqref{eq:8-coercivity}. If
\[
 \mathcal A=\{u\in g+W^{1,q}_0(U;\R^3):\det Du=1\text{ a.e.}\}
\]
has a finite-energy member, then it has an energy minimizer.
\end{theorem}
\begin{expandedproof}[weak closure of the determinant condition]
Coercivity bounds a minimizing sequence in $W^{1,q}$.
Its weak limit keeps the trace, and Lemma~\ref{lem:8-det-weak} gives
$\det Du_k\rightharpoonup\det Du$ in $L^{q/3}$.
Since every determinant equals one, the limit does too. Convex lower
semicontinuity completes existence. The feasible set is not convex,
so convexity of $L$ alone supplies no uniqueness assertion here.
The determinant condition is local and does not by itself prove global
injectivity or prevent self-interpenetration.
\end{expandedproof}
